\Subhead{2.10.1}{Самофокусировка света. Нелинейная квазиоптика и нелинейная геометрическая оптика}

В среде Керра показатель $n=n_0+n_2I$ зависит от интенсивности. Для $n_2>0$ центр пучка с большим $I$ имеет больший оптический путь и действует как собирающая линза. Параксиальная огибающая, нормированная так, что $|A|^2=I$, подчиняется
\begin{equation}
 \ii A_z+\frac1{2k}\Delta_\perp A+k_0n_2|A|^2A=0,
 \quad k=n_0k_0.
\end{equation}
Первый член после производной описывает дифракцию, второй — нелинейную фазу. Характерные длины $L_{\rm dif}\sim kw_0^2$, $L_{\rm nl}=(k_0|n_2|I_0)^{-1}$. При сравнимых длинах возможен сильный обмен между фазой и профилем; при $n_2<0$ возникает самодефокусировка.

Запись $A=\sqrt I\e^{\ii\phi}$ даёт уравнение переноса интенсивности и уравнение фазы
\begin{equation}
 \phi_z+\frac{|\nabla_\perp\phi|^2}{2k}-k_0n_2I
 =\frac1{2k}\frac{\Delta_\perp\sqrt I}{\sqrt I}.
\end{equation}
Правая часть — дифракционная поправка. Её отбрасывание даёт нелинейную геометрическую оптику, но именно она становится важной при сужении пучка. В двумерной поперечной кубической модели существует критическая мощность порядка $\lambda^2/(n_0n_2)$: выше неё возможен коллапс идеализированной модели. В реальном материале его ограничивают насыщение, высшие нелинейности, поглощение, плазма и непараксиальность. Универсального точного порога для любой формы пучка нет.
